\documentclass[11pt,a4paper]{article}
\usepackage[utf8]{inputenc}
\usepackage[T1]{fontenc}
\usepackage{geometry}
\geometry{margin=1in}
\usepackage{fancyhdr}
\usepackage{hyperref}
\usepackage{titlesec}
\usepackage{listings}
\usepackage{xcolor}

% Custom colors
\definecolor{codegreen}{rgb}{0,0.6,0}
\definecolor{codegray}{rgb}{0.5,0.5,0.5}
\definecolor{codepurple}{rgb}{0.58,0,0.82}
\definecolor{backcolour}{rgb}{0.95,0.95,0.92}

\lstdefinestyle{bash}{
    backgroundcolor=\color{backcolour},
    commentstyle=\color{codegreen},
    keywordstyle=\color{blue},
    numberstyle=\tiny\color{codegray},
    stringstyle=\color{codepurple},
    basicstyle=\ttfamily,
    breaklines=true,
    frame=single,
    numbers=left
}

\hypersetup{
    colorlinks=true,
    linkcolor=blue,
    filecolor=magenta,
    urlcolor=cyan,
}

\pagestyle{fancy}
\fancyhf{}
\rhead{\thepage}
\lhead{AI Wikipedia Documentation}

\title{\textbf{AI Wikipedia}\\\large Privacy-First AI Encyclopedia with Verification}
\author{}
\date{Version 1.1.0 -- March 2026}

\begin{document}

\maketitle
\thispagestyle{empty}

\vspace{2in}
\begin{center}
\Large A portable, privacy-first AI Wikipedia powered by Ollama with multi-source verification.
\end{center}

\newpage

\tableofcontents
\newpage

\section{Overview}

AI Wikipedia is a privacy-first application that generates articles using local Ollama LLM and verifies claims against 20+ authoritative sources including Lex.dk (Den Store Danske).

\section{Features}

\begin{itemize}
    \item \textbf{AI Article Generation} - Create articles using local Ollama LLM
    \item \textbf{AI Chat Interface} - CLI-style terminal for asking questions directly (always visible)
    \item \textbf{Multi-Source Verification} - Verify claims against 20+ authoritative sources
    \item \textbf{Lex.dk (Den Store Danske)} - Denmark's National Encyclopedia (most trusted)
    \item \textbf{Wikipedia \& WikiData} - International encyclopedia and structured data
    \item \textbf{Compressed Storage} - LZ4/Zstd compression for Wikipedia dumps
    \item \textbf{Import from Wikipedia} - Import real Wikipedia articles
    \item \textbf{Portable \& Syncable} - Export/import JSON for syncing across devices
    \item \textbf{Dark Mode} - Built-in theme toggle
\end{itemize}

\section{Quick Start}

\subsection{Prerequisites}

\begin{enumerate}
    \item \textbf{Ollama} - Install from \url{ollama.ai}
\end{enumerate}

\subsection{Installation}

\begin{lstlisting}[language=bash]
# Clone or copy the project
cd /path/to/ai-wikipedia

# Install dependencies
npm install

# Run in development mode
npm run dev
\end{lstlisting}

Open \url{http://localhost:1420} in your browser.

\section{Usage}

\subsection{AI Chat (Always Visible)}

The AI CLI chat is always visible under the search bar:
\begin{itemize}
    \item Select model: llama3.2, mistral, phi3, codellama
    \item Type questions directly and press Send or Enter
    \item Click "Save as Article" to save responses
\end{itemize}

\subsection{Generate Article with AI}

\begin{enumerate}
    \item Click "+ Generate with AI"
    \item Enter a topic (e.g., "Quantum Computing")
    \item Click Generate
    \item Article is created using Ollama locally
\end{enumerate}

\subsection{Verify Claims}

\begin{enumerate}
    \item Open any article
    \item Click "Verify Claims"
    \item Each claim is checked against enabled sources
    \item Results show: Verified, Disputed, or Unverified
\end{enumerate}

\subsection{Enable/Disable Sources}

\begin{enumerate}
    \item Click "Sources" in the header
    \item Toggle sources on/off
    \item Add custom API endpoints
    \item Changes persist in localStorage
\end{enumerate}

\subsection{Import from Wikipedia}

\begin{enumerate}
    \item Click "Import Wikipedia"
    \item Enter article titles (one per line)
    \item Click "Import"
    \item Articles are added with verified status
\end{enumerate}

\subsection{Export / Sync Data}

\begin{enumerate}
    \item Click "Export / Sync"
    \item Click "Export All" to download JSON
    \item Save to USB or cloud storage
    \item On target device, use "Import from JSON"
\end{enumerate}

\subsection{Compressed Storage}

\begin{enumerate}
    \item Click "Storage" in the header
    \item Choose compression: LZ4 (fast) or Zstd (better ratio)
    \item Click "Save to Compressed Storage"
    \item View compression stats: original vs compressed size
\end{enumerate}

\section{Verification Sources}

\subsection{Primary Sources (Enabled by Default)}

\begin{tabular}{|l|l|p{5cm}|}
    \hline
    Source & Type & Description \\
    \hline
    Lex.dk & Encyclopedia & Denmark's National Encyclopedia \\
    Wikipedia & Encyclopedia & International encyclopedia \\
    WikiData & Structured Data & Structured knowledge base \\
    \hline
\end{tabular}

\subsection{Danish Public APIs}

Available APIs: Danmarks Statistik, CVR.dk, Retsinformation, Nationalbanken, Dataforsyningen, DMI Vejr, MotorAPI, SSI COVID-19, Folketinget, DSB Feeds, Energi Data Service, elprisenligenu.dk

\section{Compression}

\subsection{Expected Ratios}

\begin{tabular}{|l|l|l|l|}
    \hline
    Algorithm & Ratio & Speed & CPU Impact \\
    \hline
    LZ4 & ~40\% & 2GB/s & Very Low \\
    Zstd & ~35\% & 500MB/s & Low-Med \\
    \hline
\end{tabular}

For 16GB Wikipedia: ~5-6GB with LZ4, ~4-5GB with Zstd.

\section{Privacy}

\begin{itemize}
    \item All AI processing uses local Ollama - no data sent to external AI services
    \item Verification uses public APIs (read-only)
    \item Data stored locally on your machine
    \item Export files are unencrypted JSON - keep them secure
\end{itemize}

\section{Tech Stack}

\begin{itemize}
    \item \textbf{Frontend}: React 18 + TypeScript + TailwindCSS
    \item \textbf{Build Tool}: Vite 5.4
    \item \textbf{Desktop}: Tauri 2.0 (Rust)
    \item \textbf{Compression}: LZ4 + Zstd (Rust)
    \item \textbf{AI}: Ollama (local LLM)
\end{itemize}

\section{Project Structure}

\begin{verbatim}
ai-wikipedia/
├── src/
│   ├── App.tsx          # Main application
│   ├── main.tsx         # React entry point
│   └── index.css        # Tailwind styles
├── src-tauri/           # Tauri desktop app
│   ├── src/
│   │   ├── lib.rs       # Rust commands
│   │   └── main.rs      # Entry point
│   └── Cargo.toml       # Rust dependencies
├── package.json
├── vite.config.ts
├── tailwind.config.js
├── README.md
├── SPEC.md
├── CHANGELOG.md
├── CONTEXT.md
└── HOWTHISWASBUILT.md
\end{verbatim}

\section{Troubleshooting}

\subsection{"Ollama not available"}

\begin{lstlisting}[language=bash]
# Start Ollama
ollama serve

# Check if running
curl http://localhost:11434/api/tags

# List available models
ollama list
\end{lstlisting}

\section{License}

MIT License

\end{document}
